\documentclass[main.tex]{subfiles}
\begin{document}

\section*{Introduction}

The relationship between urbanization and economic development has been widely studied \citep{ritchie2018urbanization}. Infrastructure, transport, and institutions play a crucial role in steering this nexus \citep{rodrigue2020geography}.

Currently, the Global South exhibits the highest urbanization rates globally. Due to systemic underinvestment in   public transport infrastructure, car dependency, ownership, and use are simultaneously rising, leading to high levels of traffic congestion. Traffic congestion has high opportunity costs, including increased transaction costs and productive hours spent driving in traffic. These high congestion levels, in turn, increase emissions and other negative externalities, such as reduced social well-being and worsened socioeconomic inequality \citep{Wang2019-ge}. Also, they are negatively affecting many sectors of these developing economies through distorted freight supply chains \citep{jabbarpour2018applications, miskolczi2021urban} and can spill over beyond individuals to city-wide stagnation and thus hinder national development. As such, the issue of congestion has a high priority on the political agenda in theses cities \citep{costa2017urban, miskolczi2021urban, perveen2017evaluating}. 

The urban mobility crisis is acutely felt in large African metropolises, and Cape Town, South Africa, is no exception. The city consistently ranks among the most congested metropolitan areas globally. According to the 2023 INRIX Global Traffic Scorecard, Cape Town ranked 9th globally out of 100 urban areas, with drivers losing an average of 83 extra hours in traffic annually \citep{inrix2023scorecard}.  

Solutions to reduce congestion are varied, and require systems thinking to mitigate congestion, whilst maintaining equitable levels of accessibility for all. One potentially useful transport system innovation that can help developing cities address congestion is Demand-Responsive Transport (DRT), a technology-enabled shared mobility service. DRT is typically seen as a helpful addition to developed country transport systems \citep{Ben-dor, engelhardt2019quantifying, schluter2021impact, Kaddoura.2020,maggi2021price, miskolczi2021urban, zwick2021ride, article}. The rationale behind the promising effect of DRT on congestion alleviation is its ability to increase the average occupancy rate of vehicles by dynamically consolidating multiple individual trips into fewer shared vehicle journeys. This mechanism directly reduces the vehicle miles traveled (VMT) compared to unshared private car trips or conventional ride-hailing services, especially during peak hours. Unlike traditional public transport, a DRT operation is not bound to fixed routes or predetermined schedules. Instead, its core characteristics revolve around flexibility and efficiency \citep{ronald2015simulating}. This efficiency stems from dynamic routing and scheduling, where vehicles do not operate over a fixed route or on a fixed schedule, allowing the service to adapt in real-time to demand. Furthermore, DRT utilizes shared journeys, typically involving dispatching a single vehicle to pick up multiple passengers at different points and transporting them to various destinations, often interrupting the route to accommodate new requests. 

To this day, however, the large-scale introduction of DRT fleets into real-world traffic dynamics of the Global South remains a purely theoretical consideration. Hence, the expected benefits of its adoption rest on the assumption that travellers will respond rationally to improved DRT availability and attractice pricing schemes. In practice, mode-choice decisions in Cape Town – and many Global South cities – are to a substantial degree shaped by non-monetary factors such as distrust in shared systems, known as ‘sharing anxiety’, and a strong preference for the perceived autonomy and security of private cars, even when stuck in congestion \citep{Dolins2021, Lavieri.2019}. While these socio-behavioural barriers are difficult to quantify, simulating a DRT system’s physical and economic performance remains a crucial first step to benchmark the theoretical potential of DRT. Agent-based models (ABM), for example, can be constructed to experiment with and compare accurate outcomes for policies, making it the ideal tool to capture the decentralized, dynamic interactions between individual user choices (demand) and fleet operations (supply) necessary to accurately test DRT feasibility and its system-wide impact before costly deployment. 

Against this backdrop, this paper investigates the potential effects of introducing a large-scale DRT system in Cape Town, South Africa, on congestion dynamics. Utilizing the Multi-Agent Transport Simulation (MATSim) framework, we rigorously test the sensitivity of the system through scenarios focused on critical operational and economic parameters, including varying fleet sizes, pricing strategies, and service designs. We hypothesize that DRT has the potential to significantly reduce congestion due to a decrease in private car use. The remainder of this article is structured as follows: first, we present a literature overview and, subsequently, lay out the model setup and parametrization of our scenario designs. We then present simulation results, followed by a discussion of the policy implications, before outlining avenues for future research.


\end{document}